% chapters/ch18-modelling-nulls.tex
\chapter{Modelling Nulls and Optional Data}

\begin{learningobjectives}
In this chapter you will learn the meaning of missing data.
You will distinguish unknown values from not-applicable values.
You will connect optionality in relationships to optionality in attributes.
\end{learningobjectives}

\section{Why nulls are a conceptual issue}
Nulls are often treated as a storage detail.
That is a mistake.
A null encodes meaning.
If the meaning is unclear, the data becomes ambiguous.

\section{Optionality versus missingness}
Optionality describes whether a relationship can be absent.
Missingness describes whether an attribute value is unknown or not applicable.
These are different ideas.

\subsection{Optional relationships}
An appointment must have a patient.
A patient may have zero appointments.
The relationship is optional from the patient side.

\subsection{Optional attributes}
A return date may not be present.
This does not mean the loan is optional.
It means the value is not yet known.

\begin{examplebox}
\textbf{Optional relationship vs optional attribute}\
A patient may have zero appointments.\\
A loan may have a return date after it is returned.\\
These are different kinds of optionality.
\end{examplebox}

\section{Unknown versus not applicable}
Unknown means the value exists but is not yet known.
Not applicable means the value does not exist for that instance.
This distinction should be explicit.

\subsection{Unknown}
“Unknown” is a temporary state.
It should usually be resolved over time.

\subsection{Not applicable}
“Not applicable” is permanent for that instance.
It should not be forced into a placeholder value.

\begin{examplebox}
\textbf{Unknown vs not applicable}\
A patient’s insurance number may be unknown at registration.\\
A walk-in patient may not have an insurance number at all.
\end{examplebox}

\section{Temporal missingness}
Some values are missing because the event has not occurred.
A return date before return is unknown, not applicable.
A graduation date before graduation is unknown.
After graduation, it is known.

\section{Defaults and placeholders}
Defaults can hide missingness.
Placeholders are often worse.
They make meaning ambiguous.
Use explicit rules instead.

\begin{principlebox}
\textbf{Principle}\
Every null must have a meaning.
\end{principlebox}

\section{Impact on constraints}
Constraints must match the intended meaning.
If a value can be unknown, allow it explicitly.
If a value is not applicable, model the condition that makes it inapplicable.

\subsection{Conceptual constraints}
State whether the value is required, optional, or conditional.
If conditional, write the rule in words.

\subsection{Relational mapping}
When you later map to relations, nullability follows these rules.
Do not decide nullability based on convenience.

\begin{chaptersummary}
Nulls are not just storage markers.
They encode meaning about unknown or not-applicable values.
Optionality in relationships is not the same as missingness in attributes.
Clear rules prevent silent ambiguity.
\end{chaptersummary}

\begin{exercisebox}
Consider a domain such as a clinic or library.
\begin{enumerate}[label=\arabic*.]
  \item List five attributes that might be missing.
  \item For each, classify as unknown or not applicable.
  \item Write a rule that explains when each value becomes known.
  \item Identify one relationship that is optional and explain why.
  \item Rewrite one ambiguous rule to remove null confusion.
\end{enumerate}
\end{exercisebox}
